\documentclass[11pt,letterpaper]{article}
\usepackage[utf8]{inputenc}
\usepackage[T1]{fontenc}
\usepackage[french]{babel}
\usepackage{amsmath,amssymb}
\usepackage{geometry}
\usepackage{graphicx}
\usepackage{enumitem}
\usepackage{tikz}
\usepackage{listings}
\usepackage{xcolor}
\usepackage{array}
\usepackage{booktabs}
\usepackage{fancyhdr}
\usepackage{lastpage}

\geometry{margin=2.5cm}

% Code listing style
\lstdefinestyle{java}{
    language=Java,
    basicstyle=\ttfamily\small,
    keywordstyle=\color{blue}\bfseries,
    commentstyle=\color{gray}\itshape,
    stringstyle=\color{red},
    showstringspaces=false,
    numbers=left,
    numberstyle=\tiny\color{gray},
    numbersep=5pt,
    frame=single,
    breaklines=true,
    tabsize=4
}

% Header/Footer
\pagestyle{fancy}
\fancyhf{}
\fancyhead[L]{IFT2015-H26 -- Examen Pratique Mi-Session}
\fancyhead[R]{Page \thepage\ / \pageref{LastPage}}
\fancyfoot[C]{\small Université de Montréal -- Département d'informatique et de recherche opérationnelle}
\renewcommand{\headrulewidth}{0.4pt}
\renewcommand{\footrulewidth}{0.4pt}

\begin{document}

% Title
\begin{center}
    {\Large \textbf{IFT2015 -- Structures de données}}\\[0.5em]
    {\large \textbf{Examen Pratique Mi-Session}}\\[0.5em]
    {\large Hiver 2026}\\[1em]
    \rule{\textwidth}{0.4pt}
\end{center}

\vspace{0.5em}

% Student info
\noindent
\begin{tabular}{@{}p{0.48\textwidth}p{0.48\textwidth}@{}}
    \textbf{Nom :} \hrulefill & \textbf{Prénom :} \hrulefill \\[1em]
    \textbf{Matricule :} \hrulefill & \textbf{Section :} \hrulefill \\
\end{tabular}

\vspace{1em}

% Instructions
\noindent\fbox{\parbox{\textwidth}{
\textbf{Directives :}
\begin{itemize}[nosep]
    \item Durée : \textbf{2 heures}
    \item Cet examen comporte \textbf{8 questions} réparties en 3 parties pour un total de \textbf{60 points}.
    \item Documents autorisés : \textbf{Aucun}
    \item Calculatrice : \textbf{Non autorisée}
    \item Justifiez vos réponses lorsque demandé.
    \item Écrivez lisiblement; les réponses illisibles ne seront pas corrigées.
\end{itemize}
}}

\vspace{1em}

% Grade table
\begin{center}
\begin{tabular}{|c|c|c|c|c|c|c|c|c|c|}
\hline
\textbf{Question} & 1 & 2 & 3 & 4 & 5 & 6 & 7 & 8 & \textbf{Total} \\
\hline
\textbf{Points} & /6 & /8 & /8 & /8 & /7 & /8 & /7 & /8 & \textbf{/60} \\
\hline
\textbf{Note} & & & & & & & & & \\
\hline
\end{tabular}
\end{center}

\vspace{1em}
\hrule
\vspace{1em}

%==============================================================================
% PARTIE 1 : THÉORIE
%==============================================================================
\section*{Partie 1 -- Théorie (30 points)}

%------------------------------------------------------------------------------
% Question 1
%------------------------------------------------------------------------------
\subsection*{Question 1 -- Vrai ou Faux (6 points)}

Pour chaque énoncé, indiquez s'il est \textbf{Vrai} ou \textbf{Faux}. Justifiez votre réponse en une ou deux phrases.

\begin{enumerate}[label=\alph*)]
    \item (1 pt) La complexité amortie O(1) de \texttt{ArrayList.add()} garantit que chaque appel individuel s'exécute en temps constant.

    \vspace{3cm}

    \item (1 pt) Dans une \texttt{LinkedList} Java, l'accès à l'élément du milieu (\texttt{get(n/2)}) est optimisé pour s'exécuter en O(1).

    \vspace{3cm}

    \item (1 pt) Une \texttt{Position} dans une liste positionnelle devient invalide si on insère un nouvel élément juste avant elle.

    \vspace{3cm}

    \item (1 pt) Dans un min-heap, le deuxième plus petit élément se trouve toujours à l'index 1 du tableau.

    \vspace{3cm}

    \item (1 pt) Un graphe non-orienté avec $n$ sommets peut avoir au maximum $n(n-1)$ arêtes.

    \vspace{3cm}

    \item (1 pt) La classe \texttt{java.util.Stack} est recommandée pour implémenter une pile dans du code Java moderne.

    \vspace{3cm}
\end{enumerate}

\newpage

%------------------------------------------------------------------------------
% Question 2
%------------------------------------------------------------------------------
\subsection*{Question 2 -- Tableau de complexité (8 points)}

Complétez le tableau suivant en indiquant la complexité temporelle de chaque opération.

\textbf{Note :} Pour la liste positionnelle, on suppose que la \texttt{Position} est déjà connue.

\vspace{0.5em}

\begin{center}
\renewcommand{\arraystretch}{1.8}
\begin{tabular}{|l|c|c|c|}
\hline
\textbf{Opération} & \textbf{ArrayList} & \textbf{LinkedList} & \textbf{Liste Positionnelle*} \\
\hline
\texttt{get(k)} -- accès par index & & & \\
\hline
\texttt{add(0, e)} -- insertion au début & & & \\
\hline
\texttt{add(n/2, e)} -- insertion au milieu & & & \\
\hline
\texttt{remove(position)} -- suppression & & & \\
\hline
\end{tabular}
\end{center}

\vspace{1em}

Complétez également ce tableau pour les files à priorité (heap) :

\vspace{0.5em}

\begin{center}
\renewcommand{\arraystretch}{1.8}
\begin{tabular}{|l|c|c|}
\hline
\textbf{Opération} & \textbf{Heap simple} & \textbf{Heap adaptable**} \\
\hline
\texttt{insert(k, v)} & & \\
\hline
\texttt{removeMin()} & & \\
\hline
\texttt{min()} & & \\
\hline
\texttt{replaceKey(entry, newKey)} & & \\
\hline
\end{tabular}
\end{center}

\vspace{0.5em}
\small{*Position déjà obtenue au préalable \quad **Avec structure de localisation}

\vspace{3cm}

%------------------------------------------------------------------------------
% Question 3
%------------------------------------------------------------------------------
\subsection*{Question 3 -- Choix multiples (8 points)}

Pour chaque question, \textbf{encerclez} la ou les bonnes réponses. Il peut y avoir plusieurs réponses correctes.

\vspace{0.5em}

\textbf{3.1} (2 pts) Quelles affirmations sont \textbf{vraies} concernant les piles et files ?

\begin{enumerate}[label=\Alph*)]
    \item \texttt{ArrayDeque} est plus performant que \texttt{java.util.Stack} pour implémenter une pile.
    \item Une file (Queue) permet l'accès au premier \textbf{et} au dernier élément en O(1).
    \item Une Deque peut être utilisée comme pile \textbf{ou} comme file.
    \item \texttt{java.util.Stack} hérite de \texttt{ArrayList}.
\end{enumerate}

\vspace{1.5cm}

\textbf{3.2} (2 pts) Pour un pattern producteur-consommateur avec buffer borné où le consommateur doit bloquer si le buffer est vide, quelle structure Java est la plus appropriée ?

\begin{enumerate}[label=\Alph*)]
    \item \texttt{LinkedList} avec \texttt{synchronized}
    \item \texttt{ArrayDeque}
    \item \texttt{ArrayBlockingQueue}
    \item \texttt{PriorityQueue}
\end{enumerate}

\vspace{1.5cm}

\textbf{3.3} (2 pts) Quelles affirmations sont \textbf{vraies} concernant les heaps ?

\begin{enumerate}[label=\Alph*)]
    \item Un heap binaire est toujours un arbre complet.
    \item Un heap et un arbre binaire de recherche (BST) ont la même propriété d'ordre.
    \item Dans un min-heap, le plus petit élément est toujours à la racine.
    \item La recherche d'un élément quelconque dans un heap est O(log n).
\end{enumerate}

\vspace{1.5cm}

\textbf{3.4} (2 pts) Pour un graphe de 10 000 sommets et 50 000 arêtes (non-orienté), quelle représentation utilise le moins de mémoire ?

\begin{enumerate}[label=\Alph*)]
    \item Matrice d'adjacence : environ 100 millions d'entrées
    \item Liste d'adjacence : environ 100 000 entrées
    \item Les deux utilisent la même quantité de mémoire
    \item Impossible à déterminer sans plus d'informations
\end{enumerate}

\newpage

%------------------------------------------------------------------------------
% Question 4
%------------------------------------------------------------------------------
\subsection*{Question 4 -- Questions à réponse courte (8 points)}

\textbf{4.1} (2 pts) Expliquez en quoi la liste de favoris avec heuristique \textit{move-to-front} diffère de la liste triée par fréquence. Donnez un avantage et un inconvénient de chaque approche.

\vspace{5cm}

\textbf{4.2} (2 pts) Qu'est-ce qu'un \textbf{invariant} de structure de données ? Donnez un exemple concret avec le coût associé à son maintien.

\vspace{5cm}

\textbf{4.3} (2 pts) Expliquez pourquoi une file à priorité adaptable nécessite une structure auxiliaire (comme un HashMap) pour atteindre O(log n) sur \texttt{replaceKey()}.

\vspace{5cm}

\textbf{4.4} (2 pts) Un graphe biparti peut-il contenir un cycle ? Si oui, quelle contrainte doit respecter ce cycle ?

\vspace{4cm}

\newpage

%==============================================================================
% PARTIE 2 : TRACES ET ANALYSE
%==============================================================================
\section*{Partie 2 -- Traces et Analyse (15 points)}

%------------------------------------------------------------------------------
% Question 5
%------------------------------------------------------------------------------
\subsection*{Question 5 -- Opérations sur un Heap (7 points)}

Soit un \textbf{min-heap} initialement vide.

\textbf{5.1} (4 pts) Insérez successivement les clés \texttt{7, 3, 9, 1, 5, 8, 2} dans le heap.

Montrez l'état du tableau après \textbf{chaque} insertion (en détaillant les échanges up-heap si nécessaire).

\vspace{10cm}

\textbf{5.2} (3 pts) À partir du heap final obtenu en 5.1, exécutez \textbf{deux} opérations \texttt{removeMin()} successives.

Pour chaque \texttt{removeMin()}, montrez :
\begin{itemize}
    \item L'élément retiré
    \item Les étapes du down-heap
    \item L'état final du tableau
\end{itemize}

\vspace{8cm}

\newpage

%------------------------------------------------------------------------------
% Question 6
%------------------------------------------------------------------------------
\subsection*{Question 6 -- Représentation de Graphe (8 points)}

Soit le graphe \textbf{orienté} $G$ défini par les arcs suivants :

\begin{center}
$(A \to B), (A \to C), (A \to D), (B \to C), (B \to E), (C \to D), (D \to E), (E \to A)$
\end{center}

\vspace{0.5em}

\textbf{6.1} (2 pts) Dessinez le graphe.

\vspace{5cm}

\textbf{6.2} (3 pts) Donnez la \textbf{matrice d'adjacence} de $G$ (5$\times$5).

\vspace{0.5em}
\begin{center}
\renewcommand{\arraystretch}{1.5}
\begin{tabular}{c|c|c|c|c|c|}
  & \textbf{A} & \textbf{B} & \textbf{C} & \textbf{D} & \textbf{E} \\
\hline
\textbf{A} & & & & & \\
\hline
\textbf{B} & & & & & \\
\hline
\textbf{C} & & & & & \\
\hline
\textbf{D} & & & & & \\
\hline
\textbf{E} & & & & & \\
\hline
\end{tabular}
\end{center}

\vspace{2cm}

\textbf{6.3} (2 pts) Donnez la \textbf{liste d'adjacence} de $G$.

\vspace{4cm}

\textbf{6.4} (1 pt) Complétez le tableau des degrés pour chaque sommet :

\vspace{0.5em}
\begin{center}
\renewcommand{\arraystretch}{1.5}
\begin{tabular}{|c|c|c|}
\hline
\textbf{Sommet} & \textbf{Degré sortant} & \textbf{Degré entrant} \\
\hline
A & & \\
\hline
B & & \\
\hline
C & & \\
\hline
D & & \\
\hline
E & & \\
\hline
\end{tabular}
\end{center}

\newpage

%==============================================================================
% PARTIE 3 : CODE ET CONCEPTION
%==============================================================================
\section*{Partie 3 -- Code et Conception (15 points)}

%------------------------------------------------------------------------------
% Question 7
%------------------------------------------------------------------------------
\subsection*{Question 7 -- Analyse et Correction de Code (7 points)}

Le code suivant tente de retourner les $k$ plus grands éléments d'un tableau. Il contient \textbf{au moins deux erreurs}.

\begin{lstlisting}[style=java]
public static int[] topK(int[] array, int k) {
    // Max-heap pour garder les k plus grands
    PriorityQueue<Integer> maxHeap = new PriorityQueue<>(
        Collections.reverseOrder()
    );

    for (int num : array) {
        maxHeap.offer(num);
        if (maxHeap.size() > k) {
            maxHeap.poll();  // Retirer le plus grand
        }
    }

    int[] result = new int[k];
    for (int i = 0; i < k; i++) {
        result[i] = maxHeap.poll();
    }
    return result;
}
\end{lstlisting}

\vspace{0.5em}

\textbf{7.1} (3 pts) Identifiez les erreurs et expliquez pourquoi elles sont problématiques.

\vspace{6cm}

\textbf{7.2} (2 pts) Réécrivez le code corrigé (vous pouvez écrire seulement les parties modifiées).

\vspace{5cm}

\textbf{7.3} (2 pts) Quelle est la complexité temporelle de l'algorithme corrigé ? Justifiez.

\vspace{4cm}

\newpage

%------------------------------------------------------------------------------
% Question 8
%------------------------------------------------------------------------------
\subsection*{Question 8 -- Scénario de Conception (8 points)}

\textbf{Contexte :} Vous concevez un système de gestion d'urgences hospitalières.

\textbf{Données :}
\begin{itemize}
    \item Patients arrivent avec un niveau d'urgence (1 = critique, 5 = mineur)
    \item À niveau égal, le patient arrivé en premier est traité en premier
    \item L'état d'un patient peut s'aggraver (son niveau change de 4 à 2, par exemple)
    \item ~500 patients en attente aux heures de pointe
    \item ~100 arrivées et ~100 traitements par heure
\end{itemize}

\vspace{0.5em}

\textbf{8.1} (2 pts) Une simple \texttt{Queue} FIFO suffit-elle ? Justifiez.

\vspace{4cm}

\textbf{8.2} (3 pts) Proposez une structure de données (ou combinaison de structures) qui répond aux besoins. Décrivez chaque opération principale et sa complexité :

\begin{itemize}
    \item \textbf{Arrivée} d'un patient
    \item \textbf{Traitement} du prochain patient (retirer le plus urgent)
    \item \textbf{Changement} de niveau d'urgence d'un patient
\end{itemize}

\vspace{8cm}

\textbf{8.3} (3 pts) Comment gérer le cas ``à niveau égal, premier arrivé premier traité'' ? Expliquez comment modifier votre structure ou clé de priorité.

\vspace{5cm}

\newpage

%==============================================================================
% CORRIGÉ
%==============================================================================
\section*{Corrigé}

\subsection*{Question 1 -- Vrai ou Faux}

\begin{enumerate}[label=\alph*)]
    \item \textbf{Faux.} La complexité amortie O(1) signifie que sur $n$ opérations, le coût total est O($n$). Mais certaines opérations individuelles (redimensionnement) peuvent prendre O($n$).

    \item \textbf{Faux.} Java optimise en partant du début ou de la fin selon l'index, mais il n'y a pas d'accès direct au milieu. \texttt{get(n/2)} reste O($n$) dans le pire cas.

    \item \textbf{Faux.} Une Position reste valide après insertions/suppressions d'autres éléments. Elle devient invalide uniquement quand son propre élément est supprimé.

    \item \textbf{Faux.} Le 2e plus petit est à l'index 1 \textbf{ou} 2 (un des deux enfants de la racine). L'invariant ne garantit pas l'ordre entre les enfants.

    \item \textbf{Faux.} Maximum = $n(n-1)/2$ pour un graphe non-orienté simple. La formule $n(n-1)$ est pour un graphe orienté.

    \item \textbf{Faux.} \texttt{java.util.Stack} est legacy, hérite de \texttt{Vector}, et est synchronisé. La recommandation est d'utiliser \texttt{ArrayDeque}.
\end{enumerate}

\subsection*{Question 2 -- Tableaux de complexité}

\begin{center}
\renewcommand{\arraystretch}{1.5}
\begin{tabular}{|l|c|c|c|}
\hline
\textbf{Opération} & \textbf{ArrayList} & \textbf{LinkedList} & \textbf{Liste Pos.*} \\
\hline
\texttt{get(k)} & O(1) & O($n$) & O($n$) \\
\hline
\texttt{add(0, e)} & O($n$) & O(1) & O(1) \\
\hline
\texttt{add(n/2, e)} & O($n$) & O($n$) & O(1)* \\
\hline
\texttt{remove(position)} & O($n$) & O(1)* & O(1) \\
\hline
\end{tabular}
\end{center}

\begin{center}
\renewcommand{\arraystretch}{1.5}
\begin{tabular}{|l|c|c|}
\hline
\textbf{Opération} & \textbf{Heap simple} & \textbf{Heap adaptable} \\
\hline
\texttt{insert(k, v)} & O(log $n$) & O(log $n$) \\
\hline
\texttt{removeMin()} & O(log $n$) & O(log $n$) \\
\hline
\texttt{min()} & O(1) & O(1) \\
\hline
\texttt{replaceKey()} & O($n$) & O(log $n$) \\
\hline
\end{tabular}
\end{center}

\subsection*{Question 3 -- Choix multiples}

\textbf{3.1} \textbf{A, C} -- ArrayDeque est plus performant, et Deque peut servir de pile ou file. (B faux: Queue n'expose que le premier; D faux: Stack hérite de Vector)

\textbf{3.2} \textbf{C} -- ArrayBlockingQueue offre put/take bloquants et capacité bornée.

\textbf{3.3} \textbf{A, C} -- Heap = arbre complet, min à la racine. (B faux: propriétés d'ordre différentes; D faux: recherche O($n$))

\textbf{3.4} \textbf{B} -- Matrice: $10000^2 = 100M$. Liste: $10000 + 2 \times 50000 = 110000 \approx 100K$. Liste gagne.

\subsection*{Question 4 -- Réponses courtes}

\textbf{4.1}
\begin{itemize}
    \item \textbf{Liste triée}: maintient l'ordre par fréquence. \texttt{getFavorites(k)} en O($k$), mais insertion coûte O($n$).
    \item \textbf{Move-to-front}: déplace en tête à chaque accès. Bénéficie de la localité temporelle, mais \texttt{getFavorites(k)} coûte O($kn$).
\end{itemize}

\textbf{4.2} Un invariant est une propriété toujours vraie sur la structure. Ex: ``liste triée'' -- coût: insertion O($n$) au lieu de O(1), bénéfice: \texttt{min()} O(1) au lieu de O($n$).

\textbf{4.3} Sans localisation, trouver une entrée dans un heap nécessite O($n$). Avec HashMap<Entry,Index>, on localise en O(1), puis up/down-heap en O(log $n$).

\textbf{4.4} Oui, un graphe biparti peut contenir des cycles. Contrainte: le cycle doit être de longueur \textbf{paire} (alterner entre les deux ensembles).

\subsection*{Question 5 -- Heap}

\textbf{5.1} Insertions:
\begin{enumerate}
    \item \texttt{insert(7)}: [7]
    \item \texttt{insert(3)}: [7,3] $\to$ 3<7 $\to$ [3,7]
    \item \texttt{insert(9)}: [3,7,9]
    \item \texttt{insert(1)}: [3,7,9,1] $\to$ 1<7 $\to$ [3,1,9,7] $\to$ 1<3 $\to$ \textbf{[1,3,9,7]}
    \item \texttt{insert(5)}: [1,3,9,7,5] $\to$ 5>3 $\to$ \textbf{[1,3,9,7,5]}
    \item \texttt{insert(8)}: [1,3,9,7,5,8] $\to$ 8<9 $\to$ \textbf{[1,3,8,7,5,9]}
    \item \texttt{insert(2)}: [1,3,8,7,5,9,2] $\to$ 2<8 $\to$ [1,3,2,7,5,9,8] $\to$ 2>1 $\to$ \textbf{[1,3,2,7,5,9,8]}
\end{enumerate}

\textbf{5.2} removeMin():
\begin{enumerate}
    \item Retirer 1, remplacer par 8: [8,3,2,7,5,9] $\to$ min(3,2)=2, 8>2 $\to$ [2,3,8,7,5,9] $\to$ min(9)=9, 8<9 $\to$ \textbf{[2,3,8,7,5,9]}
    \item Retirer 2, remplacer par 9: [9,3,8,7,5] $\to$ min(3,8)=3, 9>3 $\to$ [3,9,8,7,5] $\to$ min(7,5)=5, 9>5 $\to$ \textbf{[3,5,8,7,9]}
\end{enumerate}

\subsection*{Question 6 -- Graphe}

\textbf{6.2} Matrice:
\begin{center}
\begin{tabular}{c|ccccc}
  & A & B & C & D & E \\
\hline
A & 0 & 1 & 1 & 1 & 0 \\
B & 0 & 0 & 1 & 0 & 1 \\
C & 0 & 0 & 0 & 1 & 0 \\
D & 0 & 0 & 0 & 0 & 1 \\
E & 1 & 0 & 0 & 0 & 0 \\
\end{tabular}
\end{center}

\textbf{6.3} Liste: A$\to$[B,C,D], B$\to$[C,E], C$\to$[D], D$\to$[E], E$\to$[A]

\textbf{6.4} Degrés: A(3,1), B(2,1), C(1,2), D(1,2), E(1,2). Total out=in=8 $\checkmark$

\subsection*{Question 7 -- Code}

\textbf{7.1} Erreurs:
\begin{enumerate}
    \item \textbf{Type de heap}: Max-heap retire le plus grand, donc on perd les grands éléments! Il faut un \textbf{min-heap} de taille $k$ pour garder les $k$ plus grands.
    \item \textbf{Taille du résultat}: Si \texttt{array.length < k}, le heap aura moins de $k$ éléments, causant une exception.
\end{enumerate}

\textbf{7.2} Code corrigé:
\begin{lstlisting}[style=java]
PriorityQueue<Integer> minHeap = new PriorityQueue<>();  // Min-heap!
// ...
int[] result = new int[minHeap.size()];  // Taille reelle
\end{lstlisting}

\textbf{7.3} O($n$ log $k$): on parcourt $n$ éléments, chaque opération heap coûte O(log $k$).

\subsection*{Question 8 -- Conception}

\textbf{8.1} Non, FIFO ne tient pas compte de l'urgence. Un patient critique attendrait derrière des cas mineurs.

\textbf{8.2} \textbf{Min-heap adaptable} (priorité = niveau d'urgence) + \textbf{HashMap<PatientId, Entry>}:
\begin{itemize}
    \item Arrivée: \texttt{insert()} O(log $n$)
    \item Traitement: \texttt{removeMin()} O(log $n$)
    \item Changement: localiser via HashMap O(1), puis \texttt{replaceKey()} O(log $n$)
\end{itemize}

\textbf{8.3} Utiliser une clé composite: \texttt{(niveau, timestamp)}. Comparer d'abord par niveau, puis par timestamp si égaux. Ainsi, à niveau égal, le patient arrivé plus tôt (timestamp plus petit) est prioritaire.

\vspace{1em}

\subsection*{Barème}

\begin{center}
\begin{tabular}{|l|c|}
\hline
\textbf{Section} & \textbf{Points} \\
\hline
Q1 -- Vrai/Faux (6 $\times$ 1 pt) & 6 \\
Q2 -- Tableaux complexité (8 cases $\times$ 1 pt) & 8 \\
Q3 -- QCM (4 $\times$ 2 pts) & 8 \\
Q4 -- Réponses courtes (4 $\times$ 2 pts) & 8 \\
Q5 -- Trace heap (4 + 3 pts) & 7 \\
Q6 -- Graphe (2 + 3 + 2 + 1 pts) & 8 \\
Q7 -- Code (3 + 2 + 2 pts) & 7 \\
Q8 -- Conception (2 + 3 + 3 pts) & 8 \\
\hline
\textbf{Total} & \textbf{60} \\
\hline
\end{tabular}
\end{center}

\end{document}
